\documentclass[10pt,a4paper]{amsart}
\usepackage[margin=19mm]{geometry}
\usepackage{amsmath,amssymb,mathtools}
\usepackage[scr=rsfs,cal=euler]{mathalfa}
\usepackage{xfrac}
\usepackage[unicode,hidelinks,bookmarks=false]{hyperref}
\newcommand{\ie}{i.e.\ }
\newcommand{\cf}{cf.\ }
\newcommand{\resp}{respectively\ }
\newcommand{\Subsection}{\subsection}
\providecommand{\nobreakdash}{\nobreak\hbox{-}}
\setlength{\emergencystretch}{2em}
\multlinegap0pt
\usepackage{bm}
\usepackage{bbm}
\usepackage{dsfont}
\usepackage{xspace}
\usepackage{cancel}
\usepackage{mathtools}
\usepackage{amsthm}
\makeatletter
\@ifundefined{theorem}{\newtheorem{theorem}{Theorem}}{}
\@ifundefined{lemma}{\newtheorem{lemma}[theorem]{Lemma}}{}
\makeatother
\newcommand{\Avee}[2]{\langle#1\rangle_{#2}}
\newcommand{\BNorm}[2]{\Big\|#1\Big\|_{#2}}
\newcommand{\D}{\mathscr D}
\newcommand{\Norm}[2]{\|#1\|_{#2}}
\newcommand{\Open}{\mathrm{Open}}
\newcommand{\Rs}{\mathscr R}
\newcommand{\R}{\mathbb R}
\newcommand{\Z}{\mathbb Z}
\newcommand{\abs}[1]{|#1|}
\newcommand{\aveL}{\textit{\L}}
\newcommand{\ave}[1]{\langle #1\rangle}
\newcommand{\bNorm}[2]{\big\|#1\big\|_{#2}}
\newcommand{\one}{\mathbf{1}}
\newcommand{\vj}{\vec{\textit{\j}}}
\newcommand{\vn}{\vec n}
\newcommand{\vp}{\vec p}
\newcommand{\vq}{\vec q}
\title[Real-variable theory of matrix-weighted multi-parameter Beso]{Real-variable theory of matrix-weighted multi-parameter Besov--Triebel--Lizorkin-type spaces}
\author{Fan Bu, Yiqun Chen, Tuomas Hyt\"onen, Dachun Yang, Wen Yuan}
\date{Source: arXiv:2603.25137v1 (CC BY 4.0)}
\begin{document}
\maketitle

\noindent\emph{Source and licence.} Real-variable theory of matrix-weighted multi-parameter Besov--Triebel--Lizorkin-type spaces. Version arXiv:2603.25137v1 (CC BY 4.0), distributed under the Creative Commons Attribution 4.0 International licence, as recorded in the arXiv licence metadata for this exact version and shown on the abstract page https://arxiv.org/abs/2603.25137v1. Full rights notice: licenses/arxiv-2603.25137-CC-BY-4.0.txt.\par\smallskip

\noindent\textit{Editorial scope.} Counted unit: the Lemma whose source label is \texttt{carl >1} in the article (source0.tex lines 3347--3358, source label \texttt{carl >1}), delivered together with the article's own complete original proof of it (source0.tex lines 3360--3505, one \texttt{proof} environment of 146 lines). No mathematics is written, repaired or re-derived: statement and proof are the article's own text line for line. Required context retained uncounted: Fefferman Stein (lines 2676--2693); >1 infty at end (lines 2682--2684); FS strong (lines 2707--2721); interp (lines 3324--3333). Every definition, display and label of those lines is retained unchanged. Omitted from this package and not redistributed: 1--2675, 2685--2706, 2722--3323, 3334--3346, 3359--3359, 3506--12564. Nothing inside a retained excerpt is omitted or altered: every retained body is delivered exactly as the article's lines are written, so no omission receipt of any kind accompanies this package. Citations: the retained excerpts carry no citation command at all, so no bibliography entry of the article is transcribed and no reference is redistributed. Reference adaptation: none was needed and none was made; every reference inside the delivered unit resolves inside it, so no unresolved reference or citation is produced. Printed slips kept literal: no printed word, symbol, hypothesis or number of the counted statement or of its proof was altered. Layout and class: the article's own class is \texttt{amsart} with packages including amssymb, mathrsfs, amsmath, amsthm, amsfonts, latexsym, color, txfonts, enumerate, hyperref, anysize; the wrapper is the lane's pinned \texttt{amsart} editorial wrapper at 10pt on A4, which restates the theorem-like environments the retained text needs and adds the article's own macro definitions that the retained text needs (\texttt{\textbackslash Avee}, \texttt{\textbackslash BNorm}, \texttt{\textbackslash D}, \texttt{\textbackslash Norm}, \texttt{\textbackslash Open}, \texttt{\textbackslash R}, \texttt{\textbackslash Rs}, \texttt{\textbackslash Z}, \texttt{\textbackslash abs}, \texttt{\textbackslash ave}, \texttt{\textbackslash aveL}, \texttt{\textbackslash bNorm}, \texttt{\textbackslash one}, \texttt{\textbackslash vj}, \texttt{\textbackslash vn}, \texttt{\textbackslash vp}, \texttt{\textbackslash vq}), with extra wrapper packages (bm, bbm, dsfont, xspace, cancel, mathtools). The delivered numbering is the unit's own. Assets: no retained span contains a figure environment, a graphics-inclusion command, a PDF-page inclusion, a table or a TikZ picture; no image file is copied into this package, the package declares no asset dependency and no image file is redistributed with it. Licence: version arXiv:2603.25137v1 of this article is distributed under the Creative Commons Attribution 4.0 International licence, as recorded in the arXiv licence metadata for this exact version and shown on the abstract page https://arxiv.org/abs/2603.25137v1. A full rights notice is delivered at licenses/arxiv-2603.25137-CC-BY-4.0.txt. Characters: every retained excerpt is pure ASCII, so the delivered engine, XeLaTeX with its default Latin Modern text font, reports no missing character.
\begin{lemma}\label{Fefferman Stein}
Let $k\in\mathbb Z_+$, $l\in\mathbb N$, $(p_j)_{j=1}^{k+l}
=(\vp,\vq)\in(0,\infty]^k\times(0,\infty]^{l}$,
and $\pi\in S_{[k+l]}$.
Fix some $i\in[k+l]$ such that $\pi(i)\in[k]$,
and suppose that for some $h\in\{0,\ldots,k+l-i+1\}$
\begin{equation}\label{>1 infty at end}
(p_{\pi(j)})_{j=i}^{k+l}\in(1,\infty)^{(k+l-i+1)-h}\times\{\infty\}^{h},
\end{equation}
i.e., all these values are in $(1,\infty]$, and all infinite values, if any, are at the end of the sequence.
Then there exists a positive constant $C$,
which depends only on $\vp$, $\vq$, and $\vn$, such that,
for every $f\in(L^{\vp}\ell^{\vq})_\pi$,
\begin{equation*}
\Norm{ \mathcal M_{\pi(i)}f }{(L^{\vp}\ell^{\vq})_\pi}
\le C  \Norm{ f }{(L^{\vp}\ell^{\vq})_\pi}.
\end{equation*}
\end{lemma}

\begin{lemma}\label{FS strong}
Let $(p_i)_{i=1}^{2k}=(\vp,\vq)\in(0,\infty]^k\times(0,\infty]^{k}$ and $\pi\in S_{[2k]}$.
Let
\begin{equation}\label{mu}
\mu(\pi):=\min\{i\in[2k]:\ \pi(i)\in[k]\},
\end{equation}
and suppose that assumption \eqref{>1 infty at end} of Lemma \ref{Fefferman Stein} is satisfied with $l=k$ and $i=\mu(\pi)$.
Then there exists a positive constant $C$,
depending only on $\vp$, $\vq$, and $\vn$, such that,
for any  $f\in (L^{\vp}\ell^{\vq})_\pi$,
\begin{equation*}
\Norm{ \mathcal M_{\vn}f }{(L^{\vp}\ell^{\vq})_\pi}
\le C  \Norm{ f }{(L^{\vp}\ell^{\vq})_\pi}.
\end{equation*}
\end{lemma}

\begin{lemma}\label{interp}
Let $p_0,p_1\in(1,\infty)$ satisfy $p_0\neq p_1$,
and let $r_0,r_1,r\in[1,\infty]$.
Assume that $\theta\in(0,1)$ and $\frac 1p=\frac{1-\theta}{p_0}+\frac{\theta}{p_1}$.
Then, for any Banach space $X$,
\begin{equation*}
(L^{p_0,r_0}(X),L^{p_1,r_1}(X))_{\theta,r}=L^{p,r}(X)
\end{equation*}
with equivalent norms.
\end{lemma}

\begin{lemma}\label{carl >1}
If $p\in(1,\infty)$, $\vq\in(1,\infty)^k$, and $s\in(p,\infty)$,
then there exists a positive constant $C$ such that,
for every sequence $\{f_{\vj}\}_{\vj\in\mathbb Z^k}\in L^p\ell^{\vq}$ and
all measurable functions $\{\gamma_{\vj}\}_{\vj\in\mathbb Z^k}$,
\begin{equation}\label{carl >1eq}
\bNorm{\{\gamma_{\vj}\Avee{f_{\vj}}{\vj}\}_{\vj\in\mathbb Z^k}}{L^p\ell^{\vq}}
\leq C
\bigg[\sup_{\Omega\in\Open(\R^{\vn})}\Big\|\sup_{\vj\in\mathbb Z^k}\one_{\Omega_{\vj}}\gamma_{\vj}\Big\|_{\aveL^s(\Omega)}\bigg]
\bNorm{\{ f_{\vj}\}_{\vj\in\mathbb Z^k}}{L^p\ell^{\vq}}.
\end{equation}
\end{lemma}

\begin{proof}
We first assume that $f_{\vj}=\Avee{f_{\vj}}{\vj}$ takes a constant value $f_R$ on each $R\in\D_{\vj}(\R^{\vn})$. Let us also write $\gamma_R:=\gamma_{\vj}|_R$ for $R\in\D_{\vj}(\R^{\vn})$, although this need not be a constant. If the supremum in the claim is zero or infinity, there is nothing to prove; otherwise we may assume by scaling that it is one. We also assume that $\Norm{\{ f_{\vj}\}}{L^p\ell^{\vq}}<\infty$, for otherwise there is nothing to prove.

For every $\kappa\in\Z$, let
\begin{equation*}
\Omega_\kappa:=\Big\{x\in\R^{\vn}:\
\Norm{\{f_{\vj}(x)\}_{\vj\in\mathbb Z^k}}{\ell^{\vq}}>2^\kappa\Big\}
\end{equation*}
and
\begin{equation*}
\Rs_\kappa:=\Big\{R\in\D(\R^{\vn}):\
\abs{R\cap\Omega_\kappa}>\frac12\abs{R}\Big\}.
\end{equation*}
We claim that, for each $R\in\mathscr D(\R^{\vn})$ with $f_R\neq 0$, there is a maximal $\kappa\in\Z$ such that $R\in\Rs_\kappa$: On the one hand, for all $x\in R$, we have $\Norm{\{f_{\vj}(x)\}_{\vj\in\mathbb Z^k}}{\ell^{\vq}}\geq\abs{f_R}$; hence $R\subset\Omega_\kappa$ and thus $R\in\Rs_\kappa$ for all $\kappa\in\mathbb Z$ with $2^\kappa<\abs{f_R}$. On the other hand,
if $R\in\Rs_\kappa$, then
\begin{equation*}
\infty>\Norm{\{f_{\vj}\}_{\vj\in\mathbb Z^k}}{L^{\vp}\ell^{\vq}}
\geq 2^\kappa\abs{\Omega_\kappa}^{\frac 1p}
\geq 2^\kappa\abs{\Omega_\kappa\cap R}^{\frac 1p}
\geq 2^\kappa(2^{-1}\abs{R})^{\frac 1p},
\end{equation*}
which cannot hold for arbitrarily large $\kappa$.
Thus, $R\in\Rs_\kappa$ for all small enough $\kappa$, and $R\notin\Rs_\kappa$ for all large enough $\kappa$. This confirms the existence of a maximal $\kappa$ with $R\in\Rs_\kappa$, and hence $R\in\Rs_\kappa\setminus\Rs_{\kappa+1}$
and $R\notin\Rs_{\widetilde\kappa}\setminus\Rs_{\widetilde\kappa+1}$
for all $\widetilde\kappa\in\mathbb Z\setminus\{\kappa\}$.
By this and Minkowski's inequality, we conclude that,
for every $x\in\mathbb R^{\vec n}$,
\begin{align*}
\Norm{\{\gamma_{\vj}(x)f_{\vj}(x)\}_{\vj\in\mathbb Z^k}}{\ell^{\vq}}
&=\Norm{\{\gamma_{R}(x)f_{R}\one_R(x)\}_{R\in\D(\R^{\vn})}}{\ell^{\vq}} \\
&=\Bigg\|\Bigg\{\sum_{\kappa\in\Z}\one_{\Rs_\kappa\setminus\Rs_{\kappa+1}}(R)
\gamma_{R}(x)f_{R}\one_R(x)\Bigg\}_{R\in\D(\R^{\vn})} \Bigg\|_{\ell^{\vq}} \\
&\leq\sum_{\kappa\in\Z}\BNorm{\Big\{\one_{\Rs_\kappa\setminus\Rs_{\kappa+1}}(R)
\gamma_{R}(x)f_{R}\one_R(x)\Big\}_{R\in\D(\R^{\vn})}}{\ell^{\vq}} \\
&\leq\sum_{\kappa\in\Z}\Big[\sup_{P\in\Rs_\kappa}\abs{\gamma_P(x)}\one_P(x)\Big]
\BNorm{\Big\{\one_{\Rs_\kappa\setminus\Rs_{\kappa+1}}(R)f_{R}\one_R(x)
\Big\}_{R\in\D(\R^{\vn})}}{\ell^{\vq}}.
\end{align*}
Taking $L^p$ norms and using Minkowski's inequality again
and then using H\"older's inequality with $\frac1p=\frac1s+\frac1r$ [where, since $s\in(p,\infty)$, also $r=\frac{s}{s-p}p\in(p,\infty)$], we obtain
\begin{align*}
\Norm{\{\gamma_{\vj}f_{\vj}\}_{\vj\in\mathbb Z^k}}{L^p\ell^{\vq}}
&\leq \sum_{\kappa\in\Z}\Big\| x\mapsto \Big[\sup_{P\in\Rs_\kappa}\abs{\gamma_P(x)}\one_P(x)\Big]
\BNorm{\Big\{\one_{\Rs_\kappa\setminus\Rs_{\kappa+1}}(R)f_{R}\one_R(x)
\Big\}_{R\in\D(\R^{\vn})}}{\ell^{\vq}}\Big\|_{L^p}\\
&\leq\sum_{\kappa\in\Z}
\BNorm{x\mapsto\sup_{P\in\Rs_\kappa}\abs{\gamma_P(x)}\one_P(x)}{L^s}
\BNorm{\Big\{x\mapsto \one_{\Rs_\kappa\setminus\Rs_{\kappa+1}}(R)f_{R}\one_R(x)\Big\}_{R\in\D(\R^{\vn})}}{L^r\ell^{\vq}} \\
&=:\sum_{\kappa\in\Z}I_\kappa\times II_\kappa.
\end{align*}

To proceed, we note that every $P\in\Rs_\kappa$ satisfies $\ave{\one_{\Omega_\kappa}}_{P}>\frac12$ by definition, and hence
\begin{equation*}
P\subset\Omega_\kappa^*
:=\Big\{x\in\R^{\vn}:\ \mathcal M_{\vn}\one_{\Omega_\kappa}(x)>\frac12\Big\}.
\end{equation*}
The lower semicontinuity of $\mathcal M_{\vn}\one_{\Omega_\kappa}$ implies that $\Omega_\kappa^*$ is open
and the strong maximal inequality (Lemma \ref{FS strong}) implies that
\begin{equation}\label{Omega*}
\abs{\Omega_\kappa^*}\lesssim\abs{\Omega_\kappa}<\infty.
\end{equation}
Hence
\begin{equation*}
\begin{split}
I_\kappa
  &\leq\bigg\|x\mapsto\sup_{P\subset\Omega_\kappa^*}\abs{\gamma_P(x)}\one_P(x)\bigg\|_{L^s}
  =\bigg\|\sup_{\vj\in\mathbb Z^k}\one_{(\Omega_\kappa^*)_{\vj}} \gamma_{\vj}\bigg\|_{L^s} \\
  &\leq\abs{\Omega_\kappa^*}^{\frac 1s}\quad\text{by \eqref{carl >1eq}, where we assumed that the supremum is $1$} \\
  &\lesssim\abs{\Omega_\kappa}^{\frac 1s}\quad\text{by \eqref{Omega*}}.
\end{split}
\end{equation*}


Concerning $II_\kappa$, for $R\notin\Rs_{\kappa+1}$,
we have $\ave{\one_{\Omega_{\kappa+1}^c}}_R\geq\frac12$ by definition,
while $R\in\Rs_{\kappa}$ implies $R\subset\Omega_{\kappa}^*$ as before. Hence
\begin{align*}
II_\kappa
&\leq\BNorm{\Big\{x\mapsto \one_{\Rs_\kappa\setminus\Rs_{\kappa+1}}(R)f_{R}\cdot 2\ave{\one_{\Omega_{\kappa+1}^c}}_R\one_R(x)\Big\}_{R\in\mathscr D(\mathbb R^{\vn})}}{L^r\ell^{\vq}}  \\
&\leq \BNorm{\Big\{\Avee{2\cdot f_{\vj}\one_{\Omega_\kappa^*\setminus\Omega_{\kappa+1}}}{\vj}\Big\}_{\vj\in\mathbb Z^k}}{L^r\ell^{\vq}}
\leq \BNorm{\Big\{\mathcal M_{\vn}(2\cdot f_{\vj}\one_{\Omega_\kappa^*\setminus\Omega_{\kappa+1}}) \Big\}_{\vj\in\mathbb Z^k}}{L^r\ell^{\vq}} \\
&\lesssim\BNorm{\Big\{ f_{\vj}\one_{\Omega_\kappa^*\setminus\Omega_{\kappa+1}} \Big\}_{\vj\in\mathbb Z^k}}{L^r\ell^{\vq}}
=\BNorm{x\mapsto \one_{\Omega_\kappa^*\setminus\Omega_{\kappa+1}}(x) \Norm{\{ f_{\vj}(x) \}_{\vj\in\mathbb Z^k} }{\ell^{\vq} } }{L^r}.
\end{align*}
For $x\notin\Omega_{\kappa+1}$, we have $ \Norm{\{ f_{\vj}(x) \} }{\ell^{\vq} }\leq 2^{\kappa+1}$ by definition, and hence
\begin{equation*}
II_{\kappa}
\lesssim \BNorm{x\mapsto \one_{\Omega_\kappa^*\setminus\Omega_{\kappa+1}}(x)  2^{\kappa+1} }{L^r}
\lesssim \abs{\Omega_\kappa^*}^{\frac 1r}2^\kappa
\lesssim \abs{\Omega_\kappa}^{\frac 1r}2^\kappa,
\end{equation*}
again by \eqref{Omega*} in the last step.

Combining the obtained estimates, we deduce that
\begin{equation}\label{carl 12}
\begin{split}
  \Norm{\{\gamma_{\vj}f_{\vj}\}_{\vj\in\mathbb Z^k}}{L^p\ell^{\vq}}
  &\lesssim\sum_{\kappa\in\Z}I_\kappa\times II_\kappa
  \lesssim\sum_{\kappa\in\Z}\abs{\Omega_\kappa}^{\frac 1s}\times\abs{\Omega_\kappa}^{\frac 1r} 2^\kappa
  =\sum_{\kappa\in\Z}\abs{\Omega_\kappa}^{\frac 1p} 2^\kappa \\
  &=\sum_{\kappa\in\Z}\abs{\{x\in\R^{\vn}:\
  \Norm{\{f_{\vj}(x)\}_{\vj\in\mathbb Z^k}}{\ell^{\vq}}>2^\kappa\}}^{\frac 1p} 2^\kappa
  \sim\Norm{\{f_{\vj}\}_{\vj\in\mathbb Z^k}}{L^{p,1}\ell^{\vq}},
\end{split}
\end{equation}
where $L^{p,1}$ is the Lorentz space.

The obtained estimate deviates from that claimed in the lemma in two ways: by having the Lorentz $L^{p,1}$ norm in place of $L^p$ on the right,
and by being only established under the additional assumption that $f_{\vj}=\Avee{f_{\vj} }{\vj}$. Both shortcomings may be lifted with the help of the interpolation result stated in Lemma \ref{interp}. Given $1<p<s<\infty$ as in the assumptions of the lemma that we are proving, we may clearly choose $p_0,p_1,\theta$ as in Lemma \ref{interp} so that in addition $1<p_0<p<p_1<s<\infty$. By the strong maximal inequality (Lemma \ref{FS strong}), we find that
\begin{equation*}
\Norm{\{\Avee{f_{\vj}}{\vj}\}_{\vj\in\mathbb Z^k}}{L^{p_i}\ell^{\vq}}
\leq \Norm{\{\mathcal M_{\vn}(f_{\vj})\}_{\vj\in\mathbb Z^k}}{L^{p_i}\ell^{\vq}}
\lesssim\Norm{\{f_{\vj}\}_{\vj\in\mathbb Z^k}}{L^{p_i}\ell^{\vq}},\quad i=0,1.
\end{equation*}
Applying Lemma \ref{interp} with $X=\ell^{\vq}$, $L^{p_i,r_i}=L^{p_i,p_i}=L^{p_i}$, and $L^{p,r}=L^{p,1}$, we deduce for the bounded linear operator $\{f_{\vj}\}_{\vj\in\mathbb Z^k}\mapsto\{\Avee{f_{\vj}}{\vj}\}_{\vj\in\mathbb Z^k}:\ L^{p_i}\ell^{\vq}\to L^{p_i}\ell^{\vq}$ the estimate
\begin{equation}\label{carl 13}
\begin{split}
\Norm{\{\Avee{f_{\vj}}{\vj}\}_{\vj\in\mathbb Z^k}}{L^{p,1}\ell^{\vq}}
&\sim\Norm{\{\Avee{f_{\vj}}{\vj}\}_{\vj\in\mathbb Z^k}}{(L^{p_0}\ell^{\vq},L^{p_1}\ell^{\vq})_{\theta,1}} \\
&\lesssim\Norm{\{f_{\vj}\}_{\vj\in\mathbb Z^k}}{(L^{p_0}\ell^{\vq},L^{p_1}\ell^{\vq})_{\theta,1}}
\sim\Norm{\{f_{\vj}\}_{\vj\in\mathbb Z^k}}{L^{p,1}\ell^{\vq}}.
\end{split}
\end{equation}

Recalling that \eqref{carl 12} was proved under the assumption that $f_{\vj}=\Avee{f_{\vj}}{\vj}$, if we now apply this bound to $\Avee{f_{\vj}}{\vj}$ in place of $f_{\vj}$, and combine this with \eqref{carl 13}, we obtain
\begin{equation}\label{carl 14}
\Norm{\{\gamma_{\vj}\Avee{f_{\vj}}{\vj}\}_{\vj\in\mathbb Z^k}}{L^p\ell^{\vq}}
\lesssim \Norm{\{\Avee{f_{\vj}}{\vj}\}_{\vj\in\mathbb Z^k}}{L^{p,1}\ell^{\vq}}
\lesssim \Norm{\{f_{\vj}\}_{\vj\in\mathbb Z^k}}{L^{p,1}\ell^{\vq}}.
\end{equation}
Since \eqref{carl 14} is valid for any exponents as in the assumptions, it is also valid for each $p_i$ in place of $p$:
\begin{equation*}
\Norm{\{\gamma_{\vj}\Avee{f_{\vj}}{\vj}\}_{\vj\in\mathbb Z^k}}{L^{p_i}\ell^{\vq}}
\lesssim \Norm{\{f_{\vj}\}_{\vj\in\mathbb Z^k}}{L^{p_i,1}\ell^{\vq}},\quad i=0,1.
\end{equation*}
Applying Lemma \ref{interp} with $X=\ell^{\vq}$, $r_i\in\{p_i,1\}$, and $r=p$, we deduce for the bounded linear operator $\{f_{\vj}\}_{\vj\in\mathbb Z^k}\mapsto\{\gamma_{\vj}\Avee{f_{\vj}}{\vj}\}_{\vj\in\mathbb Z^k}:\ L^{p_i,1}\ell^{\vq}\to L^{p_i}\ell^{\vq}$, the estimate
\begin{equation*}
\begin{split}
\Norm{\{\gamma_{\vj}\Avee{f_{\vj}}{\vj}\}_{\vj\in\mathbb Z^k}}{L^{p}\ell^{\vq}}
&\sim\Norm{\{\gamma_{\vj}\Avee{f_{\vj}}{\vj}\}_{\vj\in\mathbb Z^k}}{(L^{p_0}\ell^{\vq},L^{p_1}\ell^{\vq})_{\theta,p}} \\
&\lesssim\Norm{\{\Avee{f_{\vj}}{\vj}\}_{\vj\in\mathbb Z^k}}{(L^{p_0,1}\ell^{\vq},L^{p_1,1}\ell^{\vq})_{\theta,p}}
\sim \Norm{\{f_{\vj}\}_{\vj\in\mathbb Z^k}}{L^p\ell^{\vq}},
\end{split}
\end{equation*}
which is the assertion of Lemma \ref{carl >1}.
\end{proof}

\end{document}
